\section{Ground truth mà bên tấn công buộc phải dựng}
\label{sec:ch14-ground-truth-ke-tan-cong}

\index{ground truth!dựng qua thiết kế mô hình|(}%
Mục trước đọc \tn{phép dựng giàn}{scaffolding} như bài báo của nó viết, tức là
như một phép tấn công. Mục này đọc cùng thiết lập ấy theo hướng ngược lại, và
chỗ ngược ấy là lý do chương này thuộc về Phần~IV chứ không thuộc về một phụ
lục về an toàn.

Phần~IV mắc kẹt từ chương~\ref{ch:do-do-trung-thuc} ở một chỗ: định
nghĩa~\ref{def:ch01-do-trung-thuc} hỏi lời giải thích có phản ánh đúng những
thứ mà hành vi của mô hình phụ thuộc vào hay không, và với một mô hình có sẵn
thì không ai biết những thứ ấy là gì. Không biết thì không chấm điểm được, và
mọi họ chỉ số của chương~\ref{ch:do-do-trung-thuc} tồn tại để né chỗ không biết
ấy bằng cách đo một \tn{đại lượng thay thế}{proxy} nào đó.
Chương~\ref{ch:chi-so-khong-do-do-trung-thuc} không né mà dựng: nó chọn những
nhiệm vụ mà thiết kế của nhiệm vụ ép câu trả lời phải là gì, rồi chấm điểm chỉ
số trên những nhiệm vụ ấy.

\index{ground truth!buộc phải có để tấn công}%
Thiết lập của Slack và cộng sự dựng cùng thứ đó bằng một đường khác. Mô hình
thiên lệch $f$ trong thí nghiệm không phải một mô hình bắt được ngoài đời; các
tác giả tự viết nó ra, và viết nó theo cách cực đoan nhất có thể, tức là nó dự
đoán chỉ bằng thuộc tính được bảo vệ và không dùng gì khác~\autocite{mfooling}.
Theo phương trình~\ref{eq:ch14-phep-dung-gian}, mô hình được công bố $e$ trùng
với $f$ trên toàn bộ phân phối dữ liệu thật. Vậy trên mọi điểm thật, đầu ra của
$e$ phụ thuộc vào đúng một đặc trưng, và ai cũng biết đặc trưng ấy là đặc trưng
nào. Với một điểm dữ liệu thật, câu hỏi \enquote{lời giải thích trung thực cho
$e$ ở đây phải nói gì} có một câu trả lời xác định, biết trước, và không cần
đo.

Chỗ đáng chú ý không phải là ground truth ấy tồn tại, mà là ai đã phải dựng nó
và vì sao. Trong chương~\ref{ch:chi-so-khong-do-do-trung-thuc}, người dựng là
người muốn kiểm định chỉ số, và họ phải trả giá bằng việc thu hẹp phạm vi vào
những nhiệm vụ có cấu trúc riêng. Ở đây, người dựng là bên tấn công, và không
phải vì họ muốn kiểm định gì cả: họ buộc phải biết mô hình của mình dựa vào cái
gì thì mới giấu được cái ấy. Ground truth ở đây là phụ phẩm bắt buộc của phép
tấn công.

\index{ground truth!ba đường dựng}%
Sách vì thế có ba đường dựng ground truth chứ không phải hai, và ba đường ấy
khác nhau ở chỗ cái gì bị ép. Đường thứ nhất, của
chương~\ref{ch:chi-so-khong-do-do-trung-thuc}, ép bằng thiết kế nhiệm vụ: nhiệm
vụ được chọn sao cho lời giải đúng phải đi qua một bước nhất định. Đường thứ
hai ép bằng thiết kế dữ liệu, và mục~\ref{sec:ch14-sua-bang-nhan-qua} sẽ gặp nó
ở dạng những phương trình sinh dữ liệu tự viết. Đường thứ ba, đường của mục
này, ép bằng thiết kế chính mô hình: mô hình được viết ra để chỉ phụ thuộc vào
một thứ. Cái giá của ba đường cũng khác nhau, và cái giá của đường thứ ba là
cái giá nặng nhất, nên nó phải nói ra ngay ở đây.

\index{ground truth!giới hạn của đường thứ ba}%
Kết luận rút được từ thiết lập ấy là kết luận về mô hình $e$, một mô hình được
dựng riêng để đánh bại phương pháp giải thích. Nó không phải kết luận về những
mô hình mà LIME và SHAP thường được đem đi giải thích. Nói cho chính xác: thí
nghiệm chứng minh rằng tồn tại những mô hình mà LIME và SHAP giải thích sai, và
chứng minh bằng cách trưng ra chúng. Nó không đo được LIME và SHAP sai bao
nhiêu trên một mô hình bình thường, vì trên một mô hình bình thường thì không
ai có ground truth, tức là đúng chỗ mắc kẹt ban đầu. Một phép tấn công cho ta
ground truth và lấy đi tính đại diện, đổi cái này lấy cái kia.

Ngay cả trong phạm vi hẹp ấy thì kết quả cũng không phải một con số vô nghĩa.
Một chỉ số trung thực phải thỏa những điều kiện cần trước khi bàn tới mức độ,
và mục~\ref{sec:ch08-dieu-kien-can} đã tách riêng loại phát biểu ấy khi đọc các
\tn{phép thử ngẫu nhiên hóa}{randomization test}. Thiết lập của phép dựng giàn
là một điều kiện cần cùng loại, phát biểu ở mức phương pháp thay vì ở mức chỉ
số: một phương pháp attribution chỉ hỏi mô hình ở những điểm ngoài phân phối dữ
liệu thì không phát hiện được một mô hình được dựng để hành xử khác nhau ở hai
chỗ ấy, dù mô hình thiên lệch tới đâu. Đó là một phát biểu về mọi trường hợp
chứ không về một tập dữ liệu, và nó không cần thêm thí nghiệm nào để đúng rộng
hơn.

\index{định nghĩa độ trung thực!tập can thiệp}%
Đọc kỹ hơn một bước thì thấy vì sao phát biểu ấy lại nằm gọn trong khung mà
chương~\ref{ch:giai-thich-de-lam-gi} đã dựng. Bảng~\ref{tab:ch01-bon-quan-he}
ghi rằng độ trung thực so lời giải thích với những thứ mà hành vi phụ thuộc vào
\emph{dưới một tập can thiệp đã nêu}, và cụm cuối cùng ấy mới là chỗ phép tấn
công đứng. Lời giải thích mà LIME trả về cho $e$ không sai dưới tập can thiệp
mà LIME dùng: dưới những \tn{phép nhiễu}{perturbation} ấy, $e$ thật sự chạy
theo các đặc trưng đánh lạc hướng, và LIME thuật lại đúng. Nó chỉ sai khi ta
hỏi về tập can thiệp khác, là những thay đổi nằm trong phân phối dữ liệu thật.
Phép tấn công không phá được tính trung thực dưới một tập can thiệp cố định; nó
tách hai tập can thiệp ra rồi đẩy chúng đi hai đường.

Điều đó dẫn tới hệ quả mà mục~\ref{sec:ch14-dieu-chua-giai-quyet} sẽ dùng lại,
nên nó được phát biểu ở đây một lần. Một chỉ số hay một phương pháp chỉ định
nghĩa được độ trung thực tương đối với tập can thiệp của chính nó, và tập can
thiệp ấy dựng trên hai lựa chọn tự do đầu trong ba lựa chọn của
mục~\ref{sec:ch08-vong-lap-chung}. Cho nên câu \enquote{lời giải thích này
trung thực} chưa phải một câu hoàn chỉnh chừng nào chưa nói trung thực dưới tập
can thiệp nào, và phép dựng giàn là bằng chứng rằng chỗ khuyết ấy đủ rộng để
lọt một mô hình phân biệt chủng tộc qua cửa kiểm tra.

\index{ground truth!dựng qua thiết kế mô hình|)}%
